\section{Расстояния Левенштейна и Хэмминга. Кодовое расстояние схемы. Мощность кодирования.}

Для обнаружения и исправления ошибок в каналах связи необходимо уметь измерять «различие» между словами. В теории кодирования для этого используются две основные метрики.

\subsection*{1. Расстояние Хэмминга}

\paragraph{Определение}
\textbf{Расстоянием Хэмминга} $d_H(u, v)$ между двумя словами $u$ и $v$ \textbf{одинаковой длины} называется количество позиций (разрядов), в которых соответствующие символы различны.
\textit{Пример:} $d_H(10110, 11010) = 2$ (различия во 2-м и 3-м битах).

\paragraph{Свойства}
1. Расстояние Хэмминга является метрикой (удовлетворяет неравенству треугольника).
2. Используется в блочном кодировании, где ошибки — это только \textbf{замены} (инверсии битов), но не вставки или удаления.

\subsection*{2. Расстояние Левенштейна (редакционное расстояние)}

\paragraph{Определение}
\textbf{Расстоянием Левенштейна} между двумя словами (возможно, разной длины) называется минимальное количество операций \textbf{вставки, удаления и замены} одного символа, необходимых для превращения одного слова в другое.

\paragraph{Значение}
Эта метрика более универсальна, чем расстояние Хэмминга, так как она учитывает ошибки синхронизации (когда бит пропадает или появляется лишний). Вычисляется обычно с помощью динамического программирования.

\subsection*{3. Кодовое расстояние схемы}

\paragraph{Определение}
\textbf{Кодовым расстоянием} $d_{min}$ схемы кодирования называется \textbf{минимальное} из расстояний Хэмминга между всеми парами различных кодовых слов этой схемы:
$$d = \min_{i \neq j} d_H(w_i, w_j).$$

\paragraph{Роль кодового расстояния}
Именно кодовое расстояние определяет корректирующую способность кода:
\begin{itemize}
    \item Для обнаружения $r$ ошибок необходимо, чтобы $d \ge r + 1$.
    \item Для исправления $s$ ошибок необходимо, чтобы $d \ge 2s + 1$.
\end{itemize}
\textit{Пример:} Если $d=3$, мы можем гарантированно обнаружить 2 ошибки или исправить 1 ошибку (правило «ближайшего соседа»).

\subsection*{4. Мощность кодирования}

\paragraph{Определение}
\textbf{Мощностью кода} $M$ называется количество кодовых слов в данной схеме кодирования.
В равномерном коде длины $n$ над алфавитом мощности $q$ максимально возможная мощность — $q^n$. Однако для помехоустойчивости мы используем только часть этих слов, вводя \textbf{избыточность}.

\paragraph{Параметры кода}
Код обычно характеризуется тройкой $(n, M, d)$, где:
\begin{itemize}
    \item $n$ — длина кодового слова.
    \item $M$ — мощность (число разрешенных комбинаций).
    \item $d$ — кодовое расстояние.
\end{itemize}

\paragraph{Информационная скорость}
Отношение $\frac{\log_q M}{n}$ называется скоростью кода. Она показывает, какая часть передаваемой информации является полезной. Существует фундаментальное противоречие: чем больше кодовое расстояние (надежность), тем меньше мощность кода при фиксированной длине (ниже скорость).